\documentclass[11pt,a4paper]{article}
\usepackage[T1]{fontenc}
\usepackage[utf8]{inputenc}
\usepackage[ngerman]{babel}
\usepackage{geometry}
\usepackage{enumitem}
\usepackage{amssymb}
\geometry{margin=2.2cm}
\setlength{\parindent}{0pt}
\setlength{\parskip}{0.35em}
\emergencystretch=2em
\setlist[itemize]{label=$\square$,leftmargin=1.5em,itemsep=0.25em,topsep=0.3em}
\raggedbottom
\pagestyle{plain}

\begin{document}

\begin{center}
{\Large\bfseries Checkliste fürs Betreuergespräch}\\[0.4em]
Marco Garcia Jung
\end{center}

\section*{1. Vorhaben kurz vorstellen}

Ich verbinde eine kritische Literaturanalyse zu TabPFN mit einem größeren
praktischen Teil. Darin untersuche ich, wie sich die Leistung von
TabPFN-3.5 und LightGBM bei zusätzlichen permutierten Merkmalskopien
verändert. Die eigenen Ergebnisse ordne ich anschließend in die
vorhandene Forschung ein.

\section*{2. Die wichtigsten Entscheidungen klären}

\begin{itemize}
  \item \textbf{Eigenleistung:} Genügen kritische Literatursynthese und
        empirische Anschlussuntersuchung als BA-Beitrag?
  \item Falls nicht: Welcher konkrete zusätzliche Erkenntnisgewinn
        wird erwartet?
  \item \textbf{FF1:} Darf der Literaturteil Methoden, Einsatzmöglichkeiten
        und Evaluationsgrenzen breiter behandeln?
  \item \textbf{FF2:} Passt der Robustheitsvergleich von TabPFN-3.5
        und LightGBM?
  \item \textbf{Abgrenzung:} Ist binäre Klassifikation mit zunächst
        vier numerischen Datensätzen sinnvoll?
  \item \textbf{Gewichtung:} Wie umfangreich sollen Literatur-
        und Praxisteil ungefähr sein?
  \item \textbf{Auswertung:} Genügen Leistungskurven, gepaarte Änderungen
        und Streuung, oder werden zusätzlich statistische Tests erwartet?
  \item \textbf{Titel:} Passt der aktuelle Arbeitstitel?
\end{itemize}

\section*{3. Eigenleistung klar vermitteln}

\begin{itemize}
  \item Literatur vergleichen und kritisch bewerten, nicht nur zusammenfassen.
  \item Vorhandenen Versuchsansatz selbstständig umsetzen und begründet übertragen.
  \item Eigene Ergebnisse über mehrere Datensätze erzeugen und auswerten.
  \item Ergebnisse mit den Vorarbeiten vergleichen.
  \item Keine neue Methode oder erstmalige Robustheitsuntersuchung behaupten.
\end{itemize}

\textbf{Merksatz:} Die Eigenleistung liegt in der kritischen Synthese
und der selbstständigen Durchführung und Interpretation eines
kontrollierten Anschlussversuchs.

\section*{4. Experiment verständlich erklären}

\begin{itemize}
  \item Originalspalten bleiben erhalten; zusätzliche Kopien werden
        innerhalb der Datenpartitionen durchmischt.
  \item Zusatzmerkmale: 0\,\%, 50\,\%, 100\,\% und 200\,\% der
        ursprünglichen Merkmalszahl.
  \item Beide Modelle erhalten dieselben Datenbedingungen.
  \item Zentral sind Änderungen gegenüber der jeweiligen Ausgangsleistung,
        nicht nur die Endergebnisse.
  \item Balanced Accuracy primär, Log-Loss ergänzend.
  \item Auch stabile Leistung ist ein mögliches Ergebnis.
  \item Keine kausale Erklärung von Architektur oder Vortraining
        durch diesen Versuch.
\end{itemize}

\section*{5. Literaturstand knapp einordnen}

\begin{itemize}
  \item Ursprüngliches TabPFN: Empfindlichkeit gegenüber Zusatzmerkmalen berichtet.
  \item Neuere Arbeiten: Robustheit bereits untersucht, teilweise stabile Ergebnisse.
  \item Zhao et al.: sehr ähnliches Vorgehen, unter anderem auf realen Regressionsdaten.
  \item Liu et al.: auch TabPFN-3.5, aber anderes Design mit wiederholten Messungen.
  \item Deshalb: Anschlussuntersuchung statt unberührter Forschungsfrage.
\end{itemize}

\section*{6. Offene Punkte ehrlich benennen}

\begin{itemize}
  \item Noch keine Modellläufe oder bestätigte technische Machbarkeit.
  \item Zugang, Nutzungsbedingungen und Pilot folgen nach der inhaltlichen Abstimmung.
  \item Vorläufige Datensätze: Breast Cancer, Ionosphere, QSAR und Spambase.
  \item Credit-G nur Reserve wegen kategorialer Merkmale.
  \item Identische Merkmalsvektoren insbesondere bei Spambase:
        Behandlung vor den Experimenten festlegen.
  \item Splitanteile, Wiederholungszahlen und Tuningbudget noch nicht endgültig.
\end{itemize}

\section*{7. Organisatorisches klären}

\begin{itemize}
  \item Seitenumfang und formale Vorgaben
  \item Bearbeitungszeitraum und Abgabetermin
  \item Nächster Meilenstein: Pilot, Literaturmatrix oder finaler Versuchsplan?
  \item Nächster Besprechungstermin
\end{itemize}

\section*{Ergebnis des Gesprächs festhalten}

\textbf{Akzeptierte Forschungsfragen:}\\[1.1em]
\hrule
\vspace{1em}
\textbf{Erwarteter Beitrag / nötige Änderungen:}\\[1.1em]
\hrule
\vspace{1em}
\textbf{Bewusst nicht Teil der Arbeit:}\\[1.1em]
\hrule
\vspace{1em}
\textbf{Nächster Schritt und Termin:}\\[1.1em]
\hrule

\end{document}
